\subsection{\textbf{RQ3 : What technical characteristics and limitations define
  existing Python type analysis tools?}}

The Python ecosystem includes a wide range of type analysis tools, spanning
static, dynamic, hybrid, and, more recently, LLM-based approaches.
Each category adopts a distinct technical strategy to handle Python’s highly
dynamic language characteristics, resulting in a mix of complementary strengths
and inherent limitations.
This section classifies major tools into four categories and analyzes their
technical features and constraints.

Static type checkers such as Mypy, Pyright, and Pyre are the most widely used
tools in practice.
These systems construct a type environment primarily from user-provided type
annotations and verify type consistency through constraint solving.
They support advanced features of Python’s typing system—such as Protocols and
TypedDicts—to enable structural subtyping, and they also provide partial type
inference.
This design yields fast analysis performance and reasonable precision for
large-scale codebases, making static analysis highly practical.
However, their heavy reliance on annotations means that accuracy degrades
sharply in real-world industrial code, where annotations are often missing or
incomplete.
Moreover, they cannot fully soundly model Python’s dynamic features such as
reflection, monkey patching, and dynamic attribute creation, resulting in
unavoidable false positives and false negatives.
In short, static analysis is “fast and practical but fundamentally incomplete.”

In contrast, dynamic and hybrid analysis tools leverage runtime information to
achieve higher accuracy.
Tools such as Typeguard and contract-based libraries validate arguments, return
values, and structural properties of objects at execution time, allowing them
to detect type violations soundly during actual program runs.
Tools like MonkeyType collect type hints from execution traces, enabling
high-precision, test-driven type inference.
These approaches require little to no annotation and can recover detailed type
information with high accuracy.
However, they structurally depend on test coverage: unexecuted paths cannot be
analyzed, which becomes a critical limitation given Python’s diverse dynamic
behaviors.
Runtime instrumentation also introduces performance overhead, limiting
applicability in large-scale or latency-sensitive systems.

Recent academic work explores more advanced static techniques, including
constraint-based analysis, symbolic execution, and SSA-based type analysis.
These approaches employ refined type models such as union-type lattices,
intersection types, and path-sensitive constraint solving to achieve higher
precision.
Symbolic execution allows simultaneous exploration of multiple execution paths,
mitigating the limited coverage of traditional static analysis, while SSA-based
approaches more accurately capture flow-sensitive type changes.
However, these techniques are computationally expensive and difficult to apply
to real-world codebases, especially in the presence of Python’s irregular
dynamic features such as eval, dynamic imports, and metaclasses. Scaling issues
(e.g., parameter explosion) further limit their practical deployment.

Finally, LLM-based type analysis—using models like GPT or CodeLlama—has
recently emerged as a new direction.
LLMs can infer types with high coverage even for codebases with minimal
annotations, and they can perform semantic-level reasoning by leveraging
natural-language context and code meaning.
This enables the detection of semantic type errors (e.g., API misuse, violation
of implicit assumptions) that traditional static analyzers often miss.
However, LLMs operate as black-box models; their reasoning processes are
opaque, subject to hallucination, and fundamentally unable to guarantee
soundness.
Their outputs are non-deterministic, making formal safety properties and
reproducibility—central requirements for static analysis—difficult to achieve.

In summary, Python type analysis tools exhibit complementary strengths
depending on their underlying analysis paradigm, but no single approach can
fully capture the complexity of Python’s dynamic semantics.
Static analysis is efficient yet incomplete, dynamic analysis is precise yet
coverage- and cost-constrained, and AI-based methods are flexible but lack
soundness.
Consequently, recent research trends increasingly favor hybrid approaches that
combine static, dynamic, and AI-based techniques.
This observation underscores that Python type analysis is not a problem
solvable by any single method.

\begin{tcolorbox}[colback=white,colframe=blue!60,title=Answer to RQ3,fonttitle=\bfseries]
  Python type analysis tools can be classified into static, dynamic,
  research-oriented, and LLM-based approaches. Static analysis is fast but
  struggles with dynamic features, while dynamic and hybrid methods offer
  accuracy at the cost of coverage and runtime overhead. LLM-based approaches
  enable semantic-level reasoning but lack determinism and soundness. Overall,
  no single technique is sufficient, and hybrid combinations of static, dynamic,
  and AI-based methods are becoming increasingly important.
\end{tcolorbox}
